% 第 7 章 创新点 (目标 2 页)
\section{创新点}

\subsection{三层协同的纵深防御架构}

已有方案多停留在单层防御,或规则、或评审、或离线基准。本系统给出规则层、评审层、真实环境接入层的三层协同架构,三层独立可部署、联合提供纵深。规则层拦截已知模式,评审层兜底语义可疑调用,接入层保证策略在真实链路生效。消融实验表明单规则层拦截 40.9\%,两层协同达 100\%,体现层间互补。

\subsection{针对过度代理的结构化检测维度}

已有提示注入防御工作聚焦注入内容本身,对用户只请求查看而模型自主外发的过度代理行为覆盖不足。本系统在评审层新增过度代理与意图对齐两个维度,从用户原始请求出发评估调用合法性,并以硬性约束工程化落地:非用户明确请求的高风险工具基础分不低于 6。该维度直接对应 OWASP LLM Top 10 中的 LLM08 Excessive Agency。

\subsection{封闭失败的工程安全设计}

公开的模型评审类方案在实现上普遍采用开放失败,上游异常时倾向放行,在政务数据敏感场景不可接受。本系统在基线阶段实测发现该缺陷,评审调用异常被全部放行,随后修复为封闭失败,上游异常一律阻断并升级人工审批。该修复消除评审失效即全放行的隐患,是工程级的可靠性贡献。

\subsection{真实智能体端到端验证方法}

已有方案的评估多基于离线基准或合成数据。本系统在真实智能体上加载 50 类政务技能,跑通端到端攻击触发与拦截,提供真实工具链路的拦截证据。该验证方式回应了代理能否拦截真实工具调用这一关键工程问题,验证方法本身可复用于其他智能体防护方案的评估。
